\chapter{Từ vận hành F\&B đến yêu cầu hệ thống QRTable}

Chương này chuyển bài toán SaaS POS tích hợp đặt món qua mã QR thành tập yêu cầu cụ thể cho hệ thống QRTable. Nội dung chương làm rõ hệ thống cần phục vụ những nhóm người dùng nào, mỗi nhóm thực hiện nghiệp vụ gì, dữ liệu và trạng thái nào cần được kiểm soát, cũng như những yêu cầu chất lượng nào phải được xét đến khi thiết kế và đánh giá hệ thống.

Các yêu cầu trong chương được tổng hợp từ tài liệu nghiệp vụ, ma trận phân quyền, ma trận truy vết yêu cầu–kiểm thử và các tài liệu ghi nhận phạm vi triển khai của dự án. Vì QRTable là hệ thống đa tenant, nhiều vai trò và có các luồng trạng thái liên quan đến đơn hàng, bếp và thanh toán, phần phân tích yêu cầu cần giữ ba nguyên tắc: tác nhân phải được phân biệt rõ, nguồn trạng thái đáng tin phải ổn định, và các yêu cầu phi chức năng phải gắn với tiêu chí kiểm chứng ở các chương sau.

\section{Bối cảnh vận hành và phạm vi nghiệp vụ của QRTable}

QRTable được định vị là nền tảng POS theo mô hình SaaS dành cho nhà hàng, quán ăn và quán đồ uống. Mỗi nhà hàng sử dụng hệ thống như một tenant độc lập trên cùng nền tảng. Bên trong một tenant, hệ thống hỗ trợ các hoạt động vận hành chính: quản lý thông tin nhà hàng, quản lý thực đơn, quản lý khu vực và bàn, phát hành mã QR tại bàn, tiếp nhận đơn từ khách hàng, điều phối bếp/quầy bar, xử lý thanh toán, quản lý nhân sự vận hành, theo dõi báo cáo vận hành và quản lý gói dịch vụ của tenant.

Luồng nghiệp vụ điển hình bắt đầu từ khâu khởi tạo tenant. Quản trị nền tảng tạo tenant, gán chủ nhà hàng, chọn gói dịch vụ ban đầu và khởi tạo cấu hình thanh toán. Chủ quán sau đó quản lý nhân sự, thực đơn, bàn và mã QR. Khi khách đến nhà hàng, khách quét mã QR tại bàn để tham gia một phiên gọi món. Nhiều khách cùng bàn có thể thao tác trên cùng một giỏ hàng chung. Sau khi khách gửi đơn, nhân viên xác nhận đơn trước khi món được chuyển vào màn hình bếp (KDS). Bếp hoặc bar xử lý phiếu theo trạm, cập nhật trạng thái món, trong khi nhân viên phục vụ theo dõi trạng thái đơn và xử lý hóa đơn. Khi khách yêu cầu thanh toán, hệ thống khóa phiên gọi món, tổng hợp hóa đơn, hỗ trợ thanh toán tiền mặt hoặc VietQR/SePay, sau đó chuyển bàn sang trạng thái cần dọn dẹp trước khi sẵn sàng phục vụ lượt khách tiếp theo.

Từ góc nhìn yêu cầu, QRTable là một hệ thống vận hành nhà hàng có kiểm soát trạng thái và phân quyền, vượt ra ngoài phạm vi một thực đơn điện tử đơn lẻ. Hệ thống phải đồng bộ nhiều luồng nghiệp vụ có ràng buộc trạng thái: bàn không thể vừa đang thanh toán vừa nhận thêm món; đơn chưa được xác nhận không được trừ tồn kho; khách không được nhìn thấy đơn của bàn khác; quản lý không được thực hiện các thao tác tài chính nhạy cảm dành riêng cho chủ quán; webhook thanh toán có thể đến trùng lặp nhưng không được làm hóa đơn bị đóng nhiều lần.

Về phạm vi nghiệp vụ, hệ thống tập trung vào các miền chính sau:

\begin{itemize}
  \item \textbf{Vòng đời SaaS}: quản lý tenant, gói dịch vụ, đăng ký thuê bao, hạn mức và trạng thái hoạt động của tenant.
  \item \textbf{Thực đơn, bàn và QR}: quản lý danh mục món, trạng thái món, khu vực, bàn và mã QR theo phạm vi tenant/bàn.
  \item \textbf{Phiên khách, giỏ hàng và đơn hàng}: tạo hoặc tham gia phiên gọi món, giỏ hàng chung, gửi đơn, hủy đơn chờ xử lý và yêu cầu hỗ trợ.
  \item \textbf{POS và KDS}: xác nhận đơn, điều phối món sang bếp và bar, cập nhật phiếu, phục vụ món và theo dõi trạng thái vận hành.
  \item \textbf{Quyết toán thanh toán}: tổng hợp hóa đơn, thanh toán tiền mặt, tạo VietQR, xử lý webhook SePay và hoàn tất phiên bàn.
  \item \textbf{Báo cáo vận hành và phân tích số liệu nền tảng}: tổng hợp doanh thu, đơn hàng và tình trạng bàn cho tenant; cung cấp số liệu tổng hợp cho quản trị nền tảng.
  \item \textbf{Bảo mật, RBAC và cô lập tenant}: xác thực nhân viên/quản trị bằng JWT, kiểm soát quyền theo vai trò và cô lập dữ liệu theo tenant.
\end{itemize}

\section{Bản đồ tác nhân, phạm vi truy cập và trường hợp sử dụng cốt lõi}

QRTable có hai nhóm tác nhân lớn: nhóm nội bộ vận hành nền tảng/nhà hàng và nhóm theo phiên của khách hàng. Nhóm nội bộ sử dụng tài khoản, vai trò và quyền được kiểm soát qua hệ thống xác thực và phân quyền. Khách hàng không phải là một vai trò trong cơ sở dữ liệu phân quyền; khách vào hệ thống qua mã QR và phiên của bàn. Phân biệt này quan trọng vì nếu xem khách như người dùng đăng nhập thông thường, yêu cầu bảo mật và trải nghiệm đặt món qua QR sẽ bị hiểu sai.

Hình~\ref{fig:chapter3-actor-use-case-overview} minh họa quan hệ giữa các nhóm tác nhân và các cụm trường hợp sử dụng ở mức phân tích yêu cầu bằng sơ đồ trường hợp sử dụng UML (PlantUML). Sơ đồ này cố ý tách khách theo phiên QR khỏi các vai trò RBAC của nhân viên/quản trị, đồng thời đặt cô lập tenant và quyền sở hữu phiên như yêu cầu xuyên suốt thay vì một chức năng riêng lẻ.

\begin{figure}[htbp]
  \centering
  \includegraphics[width=\textwidth]{chapter3-actor-use-case-overview.pdf}
  \caption{Tổng quan tác nhân và nhóm trường hợp sử dụng chính của QRTable.}
  \label{fig:chapter3-actor-use-case-overview}
\end{figure}

\begin{longtable}{@{}>{\raggedright\arraybackslash}p{0.18\textwidth}>{\raggedright\arraybackslash}p{0.22\textwidth}>{\raggedright\arraybackslash}p{0.54\textwidth}@{}}
\caption{Tác nhân, phạm vi truy cập và trường hợp sử dụng chính của QRTable.}
\label{tab:chapter3-actors}\\
\toprule
\textbf{Tác nhân} & \textbf{Phạm vi} & \textbf{Trường hợp sử dụng chính} \\
\midrule
\endfirsthead
\toprule
\textbf{Tác nhân} & \textbf{Phạm vi} & \textbf{Trường hợp sử dụng chính} \\
\midrule
\endhead
Quản trị nền tảng & Toàn nền tảng, chéo tenant & Đưa tenant vào hệ thống, quản lý trạng thái tenant, quản lý gói giá, xem đăng ký/hóa đơn ở lớp nền tảng, hỗ trợ xác nhận hóa đơn gói dịch vụ khi cần, xem số liệu tổng hợp của nền tảng và xem chi tiết báo cáo theo tenant qua quyền \texttt{report.read\_any}. \\
\addlinespace
Chủ nhà hàng & Một tenant & Quản lý thực đơn, danh mục, bàn, QR, nhân sự, đơn, KDS, thanh toán, gói thuê bao của tenant; thanh toán gói; cập nhật cấu hình thanh toán; xem bảng điều khiển và báo cáo của tenant qua \texttt{report.read\_own} khi gói thuê bao đang hiệu lực và gói có \texttt{analytics\_basic}. \\
\addlinespace
Quản lý & Một tenant & Vận hành nhà hàng gần giống chủ quán ở lớp nghiệp vụ hằng ngày: quản lý thực đơn/bàn/đơn/KDS/thanh toán, quản lý nhân sự trong phạm vi được cấp và xem bảng điều khiển/báo cáo của tenant; không thanh toán gói, không cập nhật cấu hình thanh toán và không xóa người dùng. \\
\addlinespace
Phục vụ/POS & Một tenant, phạm vi sàn & Xem đơn, xác nhận hoặc hủy đơn chờ xử lý, xử lý thanh toán tiền mặt, xem lịch sử thanh toán, chuyển bàn, cập nhật trạng thái bàn và xử lý yêu cầu hỗ trợ. \\
\addlinespace
Bếp/Bar & Một tenant, theo trạm bếp/bar & Xem hàng đợi KDS theo trạm, cập nhật phiếu, thu hồi phiếu thao tác nhầm; không mặc định có quyền quản lý đầy đủ thực đơn, nhân sự, thanh toán hoặc POS. \\
\addlinespace
Khách hàng & Một phiên/bàn & Quét QR, tham gia hoặc tạo phiên, xem thực đơn công khai, thao tác giỏ chung, gửi đơn, theo dõi trạng thái đơn, yêu cầu hỗ trợ, yêu cầu thanh toán và thanh toán hóa đơn hiện tại. \\
\addlinespace
Webhook SePay & Hệ thống ngoài, theo mã thanh toán & Gửi sự kiện thanh toán cho hóa đơn nhà hàng hoặc hóa đơn gói dịch vụ; không phải tác nhân người dùng và phải đi qua tuyến xác minh tương ứng. \\
\bottomrule
\end{longtable}

Các trường hợp sử dụng của quản trị nền tảng thuộc lớp nền tảng, không nằm trong tenant vận hành hằng ngày. Quản trị nền tảng có thể tạo tenant, quản lý vòng đời tenant và quản lý gói dịch vụ. Tuy nhiên, các thao tác vận hành nhà hàng thông thường vẫn thuộc chủ quán, quản lý và nhân viên trong tenant. Điều này giúp tránh tình trạng một tác nhân nền tảng bị mô tả như nhân viên vận hành của từng nhà hàng.

Chủ nhà hàng là tác nhân có phạm vi quyền rộng nhất trong tenant. Chủ quán chịu trách nhiệm thiết lập dữ liệu vận hành ban đầu như thực đơn, khu vực, bàn, mã QR, nhân sự và cấu hình thanh toán. Chủ quán cũng có quyền thực hiện các thao tác nhạy cảm với gói thuê bao SaaS và cấu hình thanh toán của tenant. Quản lý có quyền vận hành mạnh nhưng bị giới hạn ở các thao tác có rủi ro tài chính hoặc quản trị nhân sự, ví dụ không được thanh toán gói, không được cập nhật cấu hình thanh toán và không được xóa người dùng. Đây là ranh giới quan trọng để chương thiết kế sau này giải thích RBAC và chuỗi bộ kiểm soát (guard).

Nhân viên phục vụ tập trung vào luồng POS: nhận thông báo đơn chờ xử lý, xác nhận đơn, xử lý hủy đơn chờ, quản lý bàn, chuyển bàn và thanh toán. Bếp/bar tập trung vào KDS, tức xem hàng đợi và cập nhật phiếu theo trạm. Hai nhóm này đều là nhân sự của tenant nhưng có trách nhiệm khác nhau; việc tách phục vụ khỏi bếp/bar giúp yêu cầu hệ thống phản ánh đúng luồng phối hợp giữa sàn phục vụ và bếp/quầy. Theo ma trận quyền hiện tại, Waiter, Chef và Barista không mặc định có quyền xem báo cáo doanh thu hoặc bảng điều khiển báo cáo.

Khách hàng là tác nhân theo phạm vi phiên. Khách không đăng nhập bằng Keycloak, không có vai trò trong ma trận quyền và không được truy cập dữ liệu ngoài phiên/bàn của mình. Các yêu cầu của khách xoay quanh trải nghiệm thuận tiện: quét QR, xem thực đơn, cùng bàn thao tác trên giỏ hàng chung, gửi đơn, theo dõi trạng thái và thanh toán. Các kiểm soát bảo mật đối với khách vì vậy phải dựa trên token QR, guard phiên, quyền sở hữu tenant/bàn/phiên và trạng thái bàn.

\section{Năng lực hệ thống theo các miền vận hành}

Các yêu cầu chức năng được nhóm theo miền nghiệp vụ để giữ ranh giới phân tích rõ ràng. Cách nhóm này giúp tránh mô tả yêu cầu theo màn hình client hoặc theo triển khai dịch vụ quá sớm. Trong Chương 4 và Chương 5, các miền này được ánh xạ sang ranh giới dịch vụ, trách nhiệm dữ liệu và bằng chứng triển khai.

\begin{longtable}{@{}>{\raggedright\arraybackslash}p{0.10\textwidth}>{\raggedright\arraybackslash}p{0.16\textwidth}>{\raggedright\arraybackslash}p{0.43\textwidth}>{\raggedright\arraybackslash}p{0.21\textwidth}@{}}
\caption{Yêu cầu chức năng theo miền nghiệp vụ.}
\label{tab:chapter3-functional-requirements}\\
\toprule
\textbf{ID} & \textbf{Miền} & \textbf{Yêu cầu chức năng} & \textbf{Tiêu chí kiểm soát} \\
\midrule
\endfirsthead
\toprule
\textbf{ID} & \textbf{Miền} & \textbf{Yêu cầu chức năng} & \textbf{Tiêu chí kiểm soát} \\
\midrule
\endhead
FR-SAAS-01 & Đưa tenant vào hệ thống & Hệ thống phải hỗ trợ đưa tenant vào theo hướng có quản trị hỗ trợ: tạo tenant, chủ quán ban đầu, gói thuê bao mặc định và bản ghi cấu hình thanh toán trong một quy trình nhất quán. & Tạo tenant không để lại dữ liệu mồ côi khi một bước con thất bại; chủ quán được gắn đúng tenant. \\
\addlinespace
FR-SAAS-02 & Vòng đời tenant và gói thuê bao & Hệ thống phải quản lý trạng thái tenant, gói giá, gói thuê bao đang hiệu lực, hóa đơn gói và hạn mức như số bàn, số nhân viên, số đơn/ngày. & Một tenant chỉ có một gói thuê bao đang hiệu lực; tenant bị tạm ngưng bị giới hạn thao tác mới nhưng vẫn xử lý được luồng thanh toán đã phát sinh. \\
\addlinespace
FR-USER-01 & Quản lý nhân sự tenant & Chủ quán và quản lý phải quản lý nhân sự của tenant theo quyền được cấp; chủ quán có quyền thao tác nhạy hơn như xóa người dùng hoặc cập nhật vai trò, còn quản lý không được xóa người dùng. & Quản lý nhân sự thuộc phạm vi tenant, không phụ thuộc dịch vụ thông báo hoặc gửi email; không mở rộng thành quản trị nhân sự nâng cao (HRM), bảng lương hoặc lập ca. \\
\addlinespace
FR-CAT-01 & Thực đơn & Chủ quán và quản lý phải quản lý danh mục, món, ảnh, thứ tự hiển thị và trạng thái còn hàng hoặc hết hàng; khách và nhân viên chỉ xem dữ liệu thực đơn phù hợp với tenant. & Thực đơn công khai chỉ hiển thị danh mục đang hoạt động và món còn bán, chưa bị xóa. \\
\addlinespace
FR-CAT-02 & Bàn và QR & Hệ thống phải quản lý khu vực, bàn, trạng thái bàn và token QR theo phạm vi tenant/bàn để khách vào đúng phiên. & Token QR không hợp lệ hoặc không khớp tenant/bàn phải bị từ chối. \\
\addlinespace
FR-ORD-01 & Phiên và giỏ chung & Khách cùng bàn phải dùng chung phiên/giỏ; các thao tác giỏ hàng phải kiểm soát phiên bản để tránh ghi đè khi nhiều thiết bị cùng thao tác. & Cập nhật giỏ lỗi thời bị từ chối và phía client phải lấy lại ảnh chụp trạng thái từ máy chủ. \\
\addlinespace
FR-ORD-02 & Gửi đơn & Gửi đơn phải tạo đơn chờ xử lý một lần theo khóa idempotency, không tạo bản nháp đơn trong cơ sở dữ liệu và xóa giỏ sau khi gửi thành công. & Gửi lặp cùng khóa idempotency không tạo đơn trùng. \\
\addlinespace
FR-ORD-03 & Xác nhận đơn và tồn kho & Nhân viên/chủ quán/quản lý phải xác nhận đơn chờ trước khi món vào bếp; tồn kho chỉ được trừ tại thời điểm xác nhận và thuộc quyền kiểm soát của miền Thực đơn (Catalog). & Đơn không được cập nhật trực tiếp dữ liệu tồn kho của Catalog; lỗi thiếu hàng phải trả về có cấu trúc cho nhân viên. \\
\addlinespace
FR-ORD-04 & Chính sách hủy & Khách chỉ được hủy đơn chờ thuộc phiên của mình; nhân viên được hủy đơn chờ; chủ quán/quản lý được hủy đơn đang xử lý với lý do và chính sách tồn kho phù hợp. & Không cho khách hủy đơn đã vào bếp; hủy đơn phải giữ vết nghiệp vụ. \\
\addlinespace
FR-ORD-05 & Hóa đơn và yêu cầu hỗ trợ & Khách có thể yêu cầu hỗ trợ và yêu cầu thanh toán; yêu cầu hóa đơn chỉ hợp lệ khi giỏ rỗng và các đơn đang hoạt động đã ở trạng thái phục vụ xong. & Khi hóa đơn chờ thanh toán, hệ thống khóa giỏ và chuyển bàn sang trạng thái thanh toán. \\
\addlinespace
FR-KDS-01 & Bếp/KDS & Miền bếp phải nhận phiếu từ đơn đã xác nhận, tách theo trạm bếp/bar, chống trùng phiếu theo tenant, đơn và trạm. & Hàng đợi KDS hiển thị đúng trạm và không tạo phiếu trùng khi nhận sự kiện lặp. \\
\addlinespace
FR-KDS-02 & Thao tác KDS & Bếp/bar phải xem hàng đợi FIFO, cập nhật trạng thái phiếu, thu hồi thao tác nhầm, nhận cảnh báo SLA và tín hiệu thời gian thực để làm mới dữ liệu. & WebSocket chỉ là tín hiệu cập nhật; ảnh chụp trạng thái vẫn lấy từ phía máy chủ. \\
\addlinespace
FR-PAY-01 & Quyết toán thanh toán & Hệ thống phải hỗ trợ thanh toán tiền mặt và VietQR/SePay cho hóa đơn nhà hàng, dùng mã thanh toán ổn định và quy tắc làm tròn VND. & Webhook thiếu tiền không đóng hóa đơn; webhook trùng không làm thanh toán bị ghi nhận nhiều lần. \\
\addlinespace
FR-PAY-02 & Lịch sử thanh toán và đối soát & Chủ quán/quản lý xem lịch sử thanh toán theo tenant ở chế độ chỉ đọc để đối soát vận hành; điều chỉnh sau thanh toán không thuộc phạm vi luận văn. & Dữ liệu lịch sử phải được lọc theo tenant; không cho phép sửa trạng thái thanh toán từ màn hình lịch sử. \\
\addlinespace
FR-DASH-01 & Báo cáo tenant & Chủ quán và quản lý phải xem được báo cáo tenant cho doanh thu, đơn hàng, bàn, thực đơn và thanh toán ở mức tổng hợp. & Yêu cầu {\scriptsize\texttt{report.read\_own}}, gói thuê bao ở trạng thái {\scriptsize\texttt{ACTIVE}} và mã tính năng {\scriptsize\texttt{analytics\_basic}}; thành phần báo cáo bị khóa theo gói không được gọi API báo cáo. \\
\addlinespace
FR-DASH-02 & Phân tích số liệu nền tảng & Quản trị nền tảng phải xem được số liệu tổng hợp của nền tảng và xem chi tiết báo cáo tenant khi cần hỗ trợ/đối soát. & Yêu cầu {\scriptsize\texttt{report.read\_any}}; kết quả của quản trị nền tảng không bị khóa bởi gói dịch vụ của tenant được chọn. \\
\addlinespace
FR-RBAC-01 & Xác thực/RBAC & API được bảo vệ phải được kiểm soát bằng chuỗi guard, ngữ cảnh tenant và ma trận quyền; điều hướng client chỉ hỗ trợ trải nghiệm người dùng, không thay thế kiểm soát ở phía máy chủ. & API phải từ chối tác nhân thiếu quyền hoặc sai tenant ngay cả khi client bị bỏ qua. \\
\bottomrule
\end{longtable}

Nhóm yêu cầu SaaS xác định QRTable là nền tảng nhiều nhà hàng, không phải phần mềm cài riêng cho từng nhà hàng. Do đó, việc tạo tenant, gắn chủ quán, quản lý gói thuê bao và hạn mức là yêu cầu nền tảng. Trạng thái tenant cũng ảnh hưởng trực tiếp tới luồng khách: tenant bị tạm ngưng không nên tiếp nhận thao tác mới tùy tiện, nhưng các luồng thanh toán đã tạo trước đó cần được xử lý an toàn để không làm thất thoát giao dịch.

Nhóm yêu cầu nhân sự được giới hạn ở quản lý nhân sự trong tenant. Chủ quán và quản lý có thể quản lý nhân sự theo ma trận phân quyền, trong đó chủ quán có quyền thao tác nhạy hơn như xóa người dùng hoặc cập nhật vai trò, còn quản lý bị giới hạn ở các thao tác vận hành được cấp. Phạm vi này không bao gồm quy trình quản trị nhân sự nâng cao (HRM) như bảng lương, ca làm, chấm công hoặc phụ thuộc vào dịch vụ email/thông báo.

Nhóm thực đơn, bàn và QR xác định điểm vào vật lý của khách hàng. Mỗi mã QR phải gắn với đúng tenant và bàn. Thực đơn công khai cần tôn trọng trạng thái danh mục/món để khách chỉ thấy các món đang được phép bán. Các thao tác thay đổi thực đơn hoặc bàn thuộc chủ quán/quản lý, trong khi khách không có quyền quản lý thực đơn. Đây là nền cho yêu cầu cô lập tenant: cùng một hệ thống có nhiều tenant nhưng thực đơn, bàn và QR của tenant này không được lẫn với tenant khác.

Nhóm phiên, giỏ và đơn là lõi của đặt món qua QR. Một bàn có thể có nhiều khách cùng thao tác, nên giỏ hàng phải được chia sẻ và có cơ chế chống ghi đè. Gửi đơn chỉ tạo đơn chính thức khi khách gửi, còn giỏ hàng trước đó là trạng thái tạm của phiên. Xác nhận của nhân viên là điểm chuyển từ mong muốn của khách sang cam kết vận hành của nhà hàng: tồn kho được kiểm tra và trừ tại đây, đơn chuyển sang đang xử lý và phiếu được gửi sang KDS. Quy tắc này giúp tránh trừ tồn kho quá sớm khi khách chỉ đang thêm món vào giỏ.

Nhóm KDS và thanh toán hoàn tất vòng đời vận hành. KDS cần ưu tiên tính đúng đắn của hàng đợi và trạm hơn là chỉ hiển thị thời gian thực. Thanh toán cần xử lý cả giao dịch do nhân viên xác nhận trực tiếp và giao dịch do webhook từ hệ thống ngoài gửi về. Vì webhook có thể đến trễ hoặc đến nhiều lần, yêu cầu idempotency là bắt buộc trong luồng thanh toán.

Nhóm bảng điều khiển và báo cáo bổ sung khả năng đọc hiểu vận hành sau khi dữ liệu giao dịch đã phát sinh. Trong phạm vi yêu cầu của QRTable, báo cáo là mô hình dữ liệu đọc chỉ đọc do các miền sở hữu dữ liệu cung cấp, ví dụ miền Thanh toán cho doanh thu đã thanh toán, miền Đơn cho đơn/hóa đơn, miền Thực đơn cho bàn/thực đơn và miền SaaS cho số liệu nền tảng. Chương 3 vì vậy xác định nhu cầu đọc tổng hợp và quyền truy cập báo cáo, còn dịch vụ phân tích riêng, kho dữ liệu OLAP hoặc quy trình phân tích nghiệp vụ chuyên sâu thuộc nhóm mở rộng sản phẩm.

Hình~\ref{fig:chapter3-business-flow} gom các yêu cầu chức năng trên thành một luồng nghiệp vụ xuyên suốt từ lúc khách quét QR cho đến khi bàn được dọn và sẵn sàng phục vụ lượt tiếp theo. Trong Chương 3, luồng này thể hiện trình tự nghiệp vụ và điều kiện kiểm soát chính; trình tự chi tiết giữa các dịch vụ được trình bày ở Chương 5.

\begin{figure}[htbp]
  \centering
  \includegraphics[width=\textwidth,height=0.82\textheight,keepaspectratio]{chapter3-business-flow.pdf}
  \caption{Luồng nghiệp vụ từ đặt món qua QR đến KDS, thanh toán và dọn bàn.}
  \label{fig:chapter3-business-flow}
\end{figure}

\section{Ràng buộc chất lượng và tiêu chí kiểm chứng}

Các yêu cầu phi chức năng của QRTable được diễn đạt dựa trên đặc thù hệ thống SaaS nhiều tenant và các thuộc tính chất lượng thường dùng trong đánh giá hệ thống phần mềm như bảo mật, độ tin cậy, khả năng bảo trì và hiệu năng \cite{iso-iec-25010-2023}. Trong phạm vi khóa luận, các yêu cầu này đóng vai trò tiêu chí thiết kế và đánh giá có kiểm soát ở Chương 4, Chương 5 và Chương 6.

\begin{longtable}{@{}>{\raggedright\arraybackslash}p{0.16\textwidth}>{\raggedright\arraybackslash}p{0.44\textwidth}>{\raggedright\arraybackslash}p{0.34\textwidth}@{}}
\caption{Yêu cầu phi chức năng và tiêu chí đánh giá.}
\label{tab:chapter3-non-functional-requirements}\\
\toprule
\textbf{YC phi chức năng} & \textbf{Yêu cầu} & \textbf{Cách đánh giá an toàn} \\
\midrule
\endfirsthead
\toprule
\textbf{YC phi chức năng} & \textbf{Yêu cầu} & \textbf{Cách đánh giá an toàn} \\
\midrule
\endhead
Cô lập tenant & Dữ liệu, API, khóa cache, tải trọng sự kiện và tệp tải lên liên quan tenant phải luôn được ràng buộc bởi ngữ cảnh tenant. & Đánh giá bằng kiểm thử dịch vụ, kiểm thử guard, ma trận truy vết và kịch bản tenant A/B đại diện. \\
\addlinespace
Bảo mật và phân quyền & Nhân viên/chủ quán/quản trị dùng JWT/OIDC và RBAC; khách dùng xác thực theo phạm vi QR/phiên; cấu hình webhook/thanh toán không được lộ khóa bí mật hoặc token ra trình duyệt. & Đánh giá bằng ma trận quyền, chuỗi guard, hợp đồng tuyến API và kịch bản bảo mật cho QR/webhook. \\
\addlinespace
Nhất quán và idempotency & Các luồng ghi ảnh hưởng trạng thái tiền hoặc đơn phải có idempotency và nguồn trạng thái đáng tin rõ ràng: gửi đơn, xác nhận tồn kho, quyết toán webhook, đánh dấu hóa đơn đã thanh toán. & Đánh giá bằng kiểm thử máy trạng thái, kiểm thử tích hợp và kịch bản phát lại/trùng lặp. \\
\addlinespace
Phản hồi thời gian thực & Hệ thống cần phản hồi kịp thời cho POS, KDS và PWA khách, trong đó sự kiện thời gian thực đóng vai trò gợi ý hoặc làm mất hiệu lực bộ nhớ đệm. & Kiểm tra client nhận gợi ý rồi lấy lại ảnh chụp trạng thái; nguồn trạng thái nghiệp vụ cuối cùng vẫn nằm ở phía máy chủ. \\
\addlinespace
Khả năng bảo trì & Hệ thống cần có ranh giới miền rõ, hợp đồng dùng chung ổn định và cấu trúc monorepo hỗ trợ thay đổi có kiểm soát. & Đánh giá bằng trách nhiệm dịch vụ, thư viện dùng chung, hợp đồng kiến trúc và khả năng truy vết yêu cầu sang module liên quan. \\
\addlinespace
Khả năng mở rộng & Kiến trúc cần cho phép mở rộng theo ranh giới dịch vụ, bộ nhớ đệm, nhóm consumer hoặc bộ chuyển đổi thời gian thực khi cần. & Đánh giá bằng lập luận kiến trúc và bằng chứng ranh giới dịch vụ; phép đo tải lớn thuộc lớp kiểm chứng thực nghiệm với môi trường và dữ liệu riêng. \\
\addlinespace
Khả kiểm thử và minh họa tái lập được & Yêu cầu ưu tiên cao cần có ánh xạ sang kiểm thử, kịch bản minh họa hoặc tư liệu bằng chứng để người đọc kiểm tra được luận điểm. & Đánh giá bằng ma trận truy vết, kết quả kiểm thử, ảnh chụp/tư liệu minh họa và danh sách kiểm tra quá trình build. \\
\bottomrule
\end{longtable}

Cô lập tenant là yêu cầu phi chức năng trung tâm vì QRTable phục vụ nhiều nhà hàng trên cùng nền tảng. Cô lập không chỉ nằm ở truy vấn cơ sở dữ liệu mà còn ở khóa cache, phòng WebSocket, token QR, tệp tải lên, mã thanh toán và guard phân quyền. Nếu một API hoặc khóa cache bỏ qua ngữ cảnh tenant, dữ liệu của tenant này có thể bị lộ sang tenant khác. Vì vậy, mọi yêu cầu chức năng ở trên đều phải được đọc cùng yêu cầu cô lập.

Bảo mật của QRTable có hai mặt khác nhau. Với nhân viên, chủ quán, quản lý và quản trị nền tảng, hệ thống cần xác thực bằng JWT/OIDC và ủy quyền bằng RBAC. Với khách, hệ thống cần tối ưu trải nghiệm quét QR nên không yêu cầu tài khoản Keycloak; thay vào đó, quyền truy cập được giới hạn bằng quyền sở hữu QR/phiên/bàn. Với webhook thanh toán, yêu cầu bảo mật nằm ở xác minh tuyến, khóa bí mật/HMAC, mã thanh toán và idempotency. Ba nhóm cơ chế này được tách biệt khi thiết kế và đánh giá.

Nhất quán và idempotency xuất hiện ở các điểm có rủi ro cao: nhiều khách cùng sửa giỏ chung, khách gửi đơn nhiều lần do thử lại, nhân viên xác nhận đơn trong khi tồn kho thay đổi, webhook SePay gửi lặp, hoặc sự kiện thanh toán được phát lại để đồng bộ sang miền Đơn. Yêu cầu hệ thống là xác định rõ mỗi trạng thái do miền nào sở hữu, thao tác nào cần idempotency và thao tác lặp có tạo tác dụng phụ trùng lặp hay không.

Phản hồi thời gian thực giúp trải nghiệm vận hành mượt hơn trong khi nguồn trạng thái có thẩm quyền vẫn nằm ở phía máy chủ. POS, KDS và PWA khách nhận gợi ý WebSocket để biết thời điểm làm mới dữ liệu; trạng thái cuối cùng được lấy từ ảnh chụp trạng thái phía máy chủ. Cách diễn đạt này giữ rõ vai trò của WebSocket mà không phóng đại tính đúng đắn của đơn, phiếu bếp hoặc thanh toán.

\section{Các vòng đời trạng thái chi phối nghiệp vụ}

QRTable có nhiều máy trạng thái nghiệp vụ liên kết với nhau. Bàn, phiên, giỏ, đơn, phiếu KDS, hóa đơn/thanh toán và tenant đều có vòng đời riêng. Một yêu cầu quan trọng là không để một trạng thái ở miền này bị cập nhật tùy tiện bởi miền khác. Ví dụ, thanh toán hoàn tất có thể kích hoạt việc đóng hóa đơn và chuyển bàn sang trạng thái dọn dẹp, nhưng không có nghĩa miền Thanh toán tự sở hữu toàn bộ vòng đời bàn; tương tự, đơn xác nhận món có thể yêu cầu miền Thực đơn trừ tồn kho, nhưng miền Đơn không sở hữu dữ liệu tồn kho.

\begin{longtable}{@{}>{\raggedright\arraybackslash}p{0.17\textwidth}>{\raggedright\arraybackslash}p{0.35\textwidth}>{\raggedright\arraybackslash}p{0.42\textwidth}@{}}
\caption{Tóm tắt các máy trạng thái nghiệp vụ chính của QRTable.}
\label{tab:chapter3-state-machines}\\
\toprule
\textbf{Máy trạng thái} & \textbf{Trạng thái chính} & \textbf{Quy tắc chuyển trạng thái quan trọng} \\
\midrule
\endfirsthead
\toprule
\textbf{Máy trạng thái} & \textbf{Trạng thái chính} & \textbf{Quy tắc chuyển trạng thái quan trọng} \\
\midrule
\endhead
Vòng đời tenant & \texttt{ACTIVE}, \texttt{SUSPENDED}, \texttt{CLOSED} & Quản trị nền tảng quản lý vòng đời; tạm ngưng giới hạn thao tác mới nhưng vẫn cho xử lý luồng thanh toán đã phát sinh; đóng là trạng thái kết thúc hợp đồng trong phạm vi hiện tại. \\
\addlinespace
Vòng đời gói thuê bao & Hóa đơn nháp/đã phát hành, gói đang hiệu lực, gói hết hạn/thay thế & Một tenant chỉ có một gói đang hiệu lực; hóa đơn gói dùng mã tham chiếu \texttt{QRSUB}; hạn mức áp cho bàn, nhân viên và số đơn/ngày. \\
\addlinespace
Vòng đời bàn/phiên & Trống, có khách, thanh toán, dọn dẹp; phiên đang hoạt động/đã đóng & Quét QR tạo hoặc tham gia phiên khi bàn trống/có khách; thanh toán và dọn dẹp chặn đơn mới; giải phóng phiên rỗng an toàn chỉ hợp lệ khi miền Đơn chứng minh phiên rỗng. \\
\addlinespace
Vòng đời giỏ/đơn & Giỏ nháp, \texttt{PENDING}, \texttt{PROCESSING}, \texttt{READY}, \texttt{SERVED}, hoàn tất/hủy & Giỏ nháp không phải đơn lưu trong CSDL; gửi đơn tạo đơn chờ một lần; nhân viên xác nhận mới trừ tồn kho; khách chỉ hủy đơn chờ của phiên mình. \\
\addlinespace
Vòng đời phiếu KDS & Chờ, đang xử lý, sẵn sàng/xong, thu hồi & Phiếu sinh từ đơn đã xác nhận, tách theo trạm, sắp xếp theo FIFO/ưu tiên/SLA; thu hồi dùng để sửa thao tác nhầm. \\
\addlinespace
Vòng đời hóa đơn/thanh toán & \texttt{OPEN}, \texttt{PENDING\_PAYMENT}, \texttt{PAID} & Yêu cầu hóa đơn cần giỏ rỗng và đơn đã phục vụ xong; thanh toán xong đóng hóa đơn/phiên và chuyển bàn sang dọn dẹp; webhook trùng phải idempotent. \\
\bottomrule
\end{longtable}

Vòng đời bàn/phiên là cầu nối giữa thế giới vật lý của nhà hàng và trạng thái phần mềm. Khi bàn trống và khách quét QR, hệ thống tạo phiên mới và đánh dấu bàn có khách. Khi bàn có khách và phiên còn hợp lệ, khách khác cùng bàn tham gia phiên hiện tại. Khi khách yêu cầu thanh toán, bàn chuyển sang trạng thái thanh toán để khóa gọi món mới. Sau khi thanh toán hoàn tất, bàn chuyển sang dọn dẹp để nhân viên dọn bàn, rồi mới trở lại trống. Quy tắc giải phóng phiên rỗng an toàn chỉ áp dụng khi có bằng chứng rằng phiên rỗng, không có hóa đơn và không có đơn đã lưu; nó không phải chức năng mở khóa cưỡng bức tùy tiện.

Vòng đời giỏ/đơn tách rõ trạng thái trước và sau khi gửi đơn. Trước khi gửi, khách chỉ thao tác trên giỏ của phiên. Sau khi gửi, hệ thống tạo đơn chờ để nhân viên xác nhận. Tồn kho không bị trừ khi khách mới gửi, vì đơn chờ vẫn có thể bị hủy hoặc bị từ chối do thiếu hàng. Tồn kho được xử lý tại thời điểm nhân viên xác nhận, khi đơn chuyển sang đang xử lý. Cách tách này phù hợp thực tế nhà hàng: khách gửi yêu cầu, nhân viên kiểm tra rồi bếp mới bắt đầu thực hiện.

Vòng đời KDS bắt đầu sau khi đơn được xác nhận. Phiếu cần được định tuyến tới bếp hoặc bar theo trạm của món. Bếp/bar thao tác trên phiếu để chuyển trạng thái, còn phục vụ/POS theo dõi món sẵn sàng để phục vụ. KDS có thể nhận gợi ý thời gian thực để màn hình cập nhật nhanh hơn, nhưng hàng đợi phiếu và trạng thái đơn vẫn phải dựa vào ảnh chụp trạng thái phía máy chủ.

Vòng đời hóa đơn/thanh toán kết thúc phiên bàn. Hóa đơn chỉ nên vào chờ thanh toán khi không còn giỏ đang mở và các đơn cần thanh toán đã sẵn sàng. Thanh toán tiền mặt do nhân viên xác nhận trực tiếp, còn VietQR/SePay được xác nhận qua webhook. Vì thanh toán liên quan tiền và đối soát, các trường hợp thiếu tiền, trùng lặp hoặc webhook sau khi đã thanh toán phải được xử lý idempotent. Sau khi hóa đơn đã thanh toán, phiên được đóng, dữ liệu giỏ/phiên tạm được dọn và bàn chuyển sang dọn dẹp.

\section{Phạm vi yêu cầu của QRTable}

Phạm vi yêu cầu của QRTable tập trung vào các năng lực cần thiết để một đơn vị thuê bao F\&B vận hành quy trình đặt món tại bàn: quản lý tenant và gói thuê bao, thực đơn, bàn và mã QR, phiên khách, giỏ món dùng chung, đơn hàng, KDS, thanh toán, phân quyền và báo cáo vận hành mức MVP. Các yêu cầu này là cơ sở để Chương 4 thiết kế ranh giới dịch vụ, Chương 5 trình bày các luồng triển khai và Chương 6 đánh giá mức đáp ứng.

Các năng lực mở rộng như quản trị nhân sự nâng cao, CRM/loyalty, kho dữ liệu phân tích, BI/AI, ứng dụng di động gốc, hàng đợi thao tác ngoại tuyến đầy đủ, tích hợp sâu với nền tảng giao đồ ăn, hoàn tiền phức tạp hoặc nhiều ngân hàng thanh toán đồng thời được tách khỏi phạm vi yêu cầu chính. Việc đặt ranh giới này giúp phần đánh giá tập trung vào các luồng cốt lõi thay vì mở rộng sang những miền sản phẩm cần thiết kế, triển khai và kiểm chứng riêng.

Đối với hiệu năng, khả năng mở rộng và kiểm chứng nhà cung cấp thanh toán, Chương 3 xác định tiêu chí yêu cầu để Chương 4, Chương 5 và Chương 6 lần lượt thiết kế, hiện thực và phân loại bằng chứng tương ứng. Các kết luận về tải lớn, độ trễ, thông lượng hoặc callback SePay/VietQR với nhà cung cấp thật được đặt trong lớp đánh giá bằng chứng ở Chương 6. Nhờ đó, yêu cầu thanh toán vẫn bao gồm webhook, mã tham chiếu, xử lý thiếu tiền và trùng lặp, nhưng không trộn lẫn kiểm thử nội bộ với minh chứng vận hành thật.

Tóm lại, Chương 3 đã xác định tác nhân, trường hợp sử dụng, yêu cầu chức năng, yêu cầu phi chức năng và các máy trạng thái chính của QRTable. Đây là cơ sở để Chương 4 giải thích thiết kế kiến trúc và ranh giới dịch vụ, Chương 5 chứng minh quá trình triển khai, và Chương 6 đánh giá mức độ đáp ứng yêu cầu mà không phóng đại kết quả.
